% ---------------- Glosario de abreviaturas ----------------
\chapter*{Glosario de abreviaturas}
\addcontentsline{toc}{chapter}{Glosario de abreviaturas}

\begin{description}[style=sameline,leftmargin=2.2cm,labelwidth=2cm,itemsep=0.2em]
  \item[ANOVA] Análisis de la varianza (del inglés \textit{analysis of variance})
  \item[CICUAL] Comité Institucional de Cuidado y Uso de Animales de Laboratorio
  \item[CONICET] Consejo Nacional de Investigaciones Científicas y Técnicas
  \item[DAPI] 4',6-diamidino-2-fenilindol; colorante fluorescente de núcleos celulares
  \item[DCX] Doblecortina (del inglés \textit{doublecortin}); marcador de neuronas inmaduras en migración
  \item[DI] Índice de discriminación
  \item[DPPH] 2,2-difenil-1-picrilhidrazilo; radical empleado para medir actividad antioxidante
  \item[DOHaD] Orígenes del desarrollo de la salud y la enfermedad (del inglés \textit{developmental origins of health and disease})
  \item[FWER] Tasa de error por familia (del inglés \textit{family-wise error rate})
  \item[GAE] Equivalentes de ácido gálico (del inglés \textit{gallic acid equivalents})
  \item[HIP] Hipocampo
  \item[HPLC] Cromatografía líquida de alta eficacia (del inglés \textit{high-performance liquid chromatography})
  \item[IByME] Instituto de Biología y Medicina Experimental
  \item[ITBA] Instituto Tecnológico de Buenos Aires
  \item[NIH] Institutos Nacionales de Salud de los Estados Unidos (del inglés \textit{National Institutes of Health})
  \item[NeuN] Antígeno nuclear neuronal (del inglés \textit{neuronal nuclei}); marcador de neuronas maduras
  \item[NOR] Reconocimiento de objeto novedoso (del inglés \textit{novel object recognition})
  \item[OF] Prueba de \textit{open field} (campo abierto)
  \item[PBS] Solución salina tamponada con fosfato (del inglés \textit{phosphate-buffered saline})
  \item[PCA] Análisis de componentes principales (del inglés \textit{principal component analysis})
  \item[PCNA] Antígeno nuclear de células en proliferación (del inglés \textit{proliferating cell nuclear antigen})
  \item[PD] Día postnatal (del inglés \textit{postnatal day})
  \item[PFA] Paraformaldehído
  \item[p/v] Peso en volumen (unidad de concentración)
  \item[QE] Equivalentes de quercetina (del inglés \textit{quercetin equivalents})
  \item[SEM] Error estándar de la media (del inglés \textit{standard error of the mean})
  \item[SLR] Reconocimiento espontáneo de posición (del inglés \textit{spontaneous location recognition})
  \item[d-SLR] Configuración \textit{dissimilar} de la prueba SLR; los tres objetos se disponen equidistantes (a 120°), lo que impone una menor demanda de separación de patrones
  \item[s-SLR] Configuración \textit{similar} de la prueba SLR; dos de los objetos se colocan próximos entre sí (separados $\sim$50°), lo que impone una mayor demanda de separación de patrones
  \item[SNC] Sistema nervioso central
  \item[v/v] Volumen en volumen (unidad de concentración)
\end{description}
